\documentclass[10pt]{article}
\usepackage[a4paper,margin=16mm]{geometry}
\usepackage[no-math]{fontspec}
\setmainfont{Latin Modern Roman}
\usepackage{amsmath,amsfonts,amssymb,amsthm,graphicx,enumitem,xurl}
\usepackage[unicode,hidelinks]{hyperref}
\newtheorem{theorem}{Theorem}[section]
\newtheorem{lemma}[theorem]{Lemma}
\newtheorem{proposition}[theorem]{Proposition}
\newtheorem{definition}[theorem]{Definition}
\newtheorem{corollary}[theorem]{Corollary}
\newtheorem{remark}[theorem]{Remark}
\allowdisplaybreaks
\setlength\emergencystretch{3em}
  \renewenvironment{proof}[1][]{
    	\begin{trivlist}
     	\item[\hspace{\labelsep}{\em\noindent Proof#1:\/}]}
     	{{\hfill$\Box$}
    	\end{trivlist}
  }
  \DeclareMathOperator{\var}{Var}
  \renewcommand{\Pr}{\mbox{\rm Pr}}
  \newcommand{\Exp}{{\mathbb{E}}}
  \newcommand{\E}{\mathbb{E}}
  \DeclareMathOperator{\gl}{GL}
  \newcommand{\R}{\mathbb{R}} 
  \newcommand{\C}{\mathbb{C}} 
  \newcommand{\D}{\mathbb{D}} 
  \newcommand{\N}{\mathbb{N}} 
  \newcommand{\Z}{\mathbb{Z}} 
  \newcommand{\F}{\mathbb{F}} 
  \newcommand{\K}{\mathbb K} 
  \newcommand{\T}{\mathbb T} 
  \newcommand{\pmset}[1]{\{-1,1\}^{#1}} 
  \newcommand{\bset}[1]{\{0,1\}^{#1}} 
  \newcommand{\sphere}[1]{S^{#1-1}} 
  \newcommand{\ball}[1]{B_{#1}} 
  \newcommand{\Orth}[1]{O(\R^{#1})} 
  \DeclareMathOperator{\im}{im} 
  \DeclareMathOperator{\vspan}{Span} 
  \newcommand{\1}{\mathbf{1}}
  \DeclareMathOperator{\gram}{Gram}
  \DeclareMathOperator{\herm}{Herm}
  \DeclareMathOperator{\conv}{Conv}
  \newcommand{\Id}{\ensuremath{\mathop{\rm I}\nolimits}}
 \newcommand{\Zp}{\mathbb Z_p} 
  \DeclareMathOperator{\charac}{char}
  \newcommand{\one}{\mathbf{1}}
  \DeclareMathOperator{\supp}{supp}
  \DeclareMathOperator{\ent}{\mathcal N}
  \DeclareMathOperator{\diam}{diam}
  \newcommand{\st}{:\,} 
  \newcommand{\ie}{{i.e.}}  
  \newcommand{\eg}{{e.g.}} 
  \newcommand{\eps}{\varepsilon}
  \newcommand{\ip}[1]{\langle #1 \rangle}
  \newcommand{\cF}{{\mathcal F}}   
  \DeclareMathOperator{\U}{\mathcal{U}}
  \DeclareMathOperator{\sign}{sign}
  \newcommand{\poly}{\mbox{\rm poly}}
  \newcommand{\ceil}[1]{\lceil{#1}\rceil}
  \newcommand{\floor}[1]{\lfloor{#1}\rfloor}
  \DeclareMathOperator{\Tr}{Tr}
  \DeclareMathOperator{\diag}{diag}
  \DeclareMathOperator{\Diag}{Diag}
  \DeclareMathOperator{\spec}{Spec}
  \DeclareMathOperator{\Fchar}{char}
  \DeclareMathOperator{\enc}{\sf Enc}
  \DeclareMathOperator{\dec}{\sf Dec}
  \DeclareMathOperator{\Norm}{N}
  \DeclareMathOperator{\rank}{rk}
  \DeclareMathOperator{\cay}{Cay}
  \DeclareMathOperator{\opt}{OPT}
  \DeclareMathOperator{\sdp}{SDP}
  \DeclareMathOperator{\cut}{cut}
  \newcommand{\id}{I}
  \newcommand{\A}{\mathcal A}
  \newcommand{\B}{\mathcal B}
  \newcommand{\CC}{\mathcal C}
  \newcommand{\HS}{\mathcal H}
  \DeclareMathOperator{\symten}{Sym}
  \DeclareMathOperator{\arank}{arank}
  \DeclareMathOperator{\prank}{prank} 
  \DeclareMathOperator{\srank}{srank} 
  \DeclareMathOperator{\GR}{GR} 
  \DeclareMathOperator{\bias}{bias}
  \DeclareMathOperator{\Pol}{Pol}
  \DeclareMathOperator{\codim}{codim}
  \DeclareMathOperator{\smrank}{schmrank}
  \newcommand{\beq}{\begin{equation}}
  \newcommand{\eeq}{\end{equation}}
  \newcommand{\beqn}{\begin{equation*}}
  \newcommand{\eeqn}{\end{equation*}}
  \newcommand{\beqr}{\begin{eqnarray}}
  \newcommand{\eeqr}{\end{eqnarray}}
  \newcommand{\beqrn}{\begin{eqnarray*}}
  \newcommand{\eeqrn}{\end{eqnarray*}}
  \newcommand{\bmline}{\begin{multline}}
  \newcommand{\emline}{\end{multline}}
  \newcommand{\bmlinen}{\begin{multline*}}
  \newcommand{\emlinen}{\end{multline*}}
  \newcommand{\nc}{{\normalfont NC}}
    \newcommand{\ncz}{{\normalfont NC}$^0$}
  \newcommand{\qnc}{{\normalfont QNC}}
  \newcommand{\ac}{{\normalfont  AC}}
    \newcommand{\acz}{{\normalfont  AC}$^0$}
  \newcommand{\acc}{{\normalfont  ACC}}
  \newcommand{\qac}{{\normalfont  QAC}}
  \newcommand{\qacc}{{\normalfont  QACC}}
\begin{document}
\begin{center}\Large Problem 2532: Random polynomial-map restrictions\end{center}
\noindent\textbf{Authors:} Jop Bri\"et and Davi Castro-Silva\par
\noindent\textbf{Work:} Random Restrictions of High-Rank Tensors and Polynomial Maps. Discrete Analysis 2024:9, DOI 10.19086/da.124610, exact arXiv 2212.13728v2 (5 November 2024), 23 physical pages.\par
\noindent\textbf{Native source and grant:} \url{https://arxiv.org/src/2212.13728v2}, CC BY4.0, \url{https://creativecommons.org/licenses/by/4.0/}. Original PDF companion separately CC BY3.0: \url{http://creativecommons.org/licenses/by/3.0/}.\par
\noindent\textbf{Difficulty:} H1: The polynomial restriction rank need not be subadditive as a Boolean set function, so the tensor concentration theorem does not supply this arbitrary-characteristic outcome. The complete authored degree induction passes to polynomial coefficient rings and extracts mixed monomials in two successive restrictions; a separate low-influence junta assignment and Lipschitz argument boosts positive probability to near certainty. These nonroutine local constructions remain after the precisely cited Margulis--Russo and Friedgut prerequisites, and are above ordinary qualifying-exam training.\par
\noindent\textbf{Editorial selection and prerequisites:} Theorem 1.8 only, arbitrary field and degree, natural polynomial-map ranks as Definition 3.1, kappa depends only on d/sigma and threshold R may additionally depend on epsilon. No exponential tail asserted. Complete Lemma 3.3 coefficient-ring degree induction, Lemma 3.4 junta boosting and final closure are retained. The tensor statement, natural-rank definition and full concentration/coupling lemma are uncounted comparison/support, not a second selected problem. Ordinary named prerequisites: Margulis--Russo formula, Friedgut 1998 junta theorem, and its biased-measure constants as O'Donnell 2014 Section 10.3. Source Lemma 3.4 literally says sigma>0; the selected main proof applies it only at 0.9 sigma, with the whole interval inside (0,1). That literal wider helper wording is preserved; no extension outside the probability range is selected.\par
\noindent\textbf{Changes:} Independent standard article/theorem wrapper and original counter restoration; published class, original preamble and unprinted comments excluded. Printed mathematical tokens retained; no proof/formula replacement. Original contents list is navigation only, not a version pin. No live-corpus admission or publication claimed.\par\medskip\hrule\medskip
\setcounter{section}{1}\setcounter{theorem}{1}
Since our results are independent of the field considered (which can be finite or infinite), we will always denote it by $\F$ and suppress statements of the form ``let $\F$ be a field'' or ``for every field $\F$''.
We begin by considering the case of tensors.

\begin{definition}[Tensors]\label{def:tensors}
For finite sets $X_1,\dots,X_d \subset \N$, a $d$-tensor is a map $T:X_1\times\cdots\times X_d\to \F$.
We will associate with any $d$-tensor~$T$ a multilinear map $\F^{X_1}\times\cdots\times \F^{X_d}\to \F$ and an element of $\F^{X_1}\otimes \cdots\otimes \F^{X_d}$ in the obvious way, and also denote these objects by~$T$.
\end{definition}

\begin{definition}[Restriction of tensors]
For a tensor~$T$ as in Definition~\ref{def:tensors} and subsets $I_1\subseteq X_1,\dots,I_d\subseteq X_d$, denote $I_{[d]} = I_1\times \cdots\times I_d$ and write~$T_{|I_{[d]}}$ for the restriction of~$T$ to~$I_{[d]}$.
If~$T$ is viewed as an element of $\F^{X_1}\otimes \cdots\otimes \F^{X_d}$, then~$T_{|I_{[d]}}$ is simply a sub-tensor.
\end{definition}

We define $(\F^\infty)^{\otimes d}$ to be the set of $d$-tensors over~$\F$ with finite support.
Note that the tensors defined on finite sets naturally embed into this set, and that the rank functions for tensors discussed above are invariant under this embedding.
Our main result regarding tensors is then as follows:

\begin{theorem}\label{thm:random_tensor}
For every~$d\in \N$ and~$\sigma\in (0,1]$, there exist constants $C, \kappa > 0$ such that the following holds.
For every natural rank function $\rank: (\F^{\infty})^{\otimes d} \to \R_+$ and every $d$-tensor $T \in \bigotimes_{i=1}^d \F^{[n_i]}$ we have
\beqn
\Pr_{I_1\sim [n_1]_{\sigma}, \dots, I_d\sim [n_d]_{\sigma}}\big[ \rank \big(T_{|I_{[d]}}\big) \geq \kappa\cdot \rank(T)\big] \geq 1 - C e^{-\kappa \rank(T)}.
\eeqn
\end{theorem}
\setcounter{theorem}{5}
\begin{definition}[Polynomial map]\label{def:polymap}
A polynomial map is an ordered tuple $\phi(x) = \big(f_1(x),\dots,f_k(x)\big)$ of polynomials $f_1,\dots,f_k\in \F[x_1,\dots,x_n]$.
We identify with~$\phi$ a map $\F^n\to\F^k$ in the natural way.
The degree of~$\phi$ is the maximum degree of the~$f_i$.
\end{definition}

\begin{definition}[Restriction of polynomial maps]
For a polynomial map~$\phi:\F^n\to\F^k$ and a set~$I\subseteq [n]$, define the restriction~$\phi_{|I}:\F^I\to\F^k$ to be the map given by
$\phi_{|I}(y) = \phi(\bar y)$,
where~$\bar y\in \F^n$ agrees with~$y$ on the coordinates in~$I$ and is zero elsewhere.
\end{definition}

We denote the space of all polynomial maps $\phi: \F^n\to\F^k$ of degree at most~$d$ by $\Pol_{\leq d}(\F^n, \F^k)$, and write
\beqn
\Pol_{\leq d}(\F^{\infty}, \F^k) = \bigcup_{n\in \N} \Pol_{\leq d}(\F^n, \F^k).
\eeqn
Our main result in this setting is the following:

\begin{theorem} \label{thm:random_poly}
For every~$d\in \N$ and~$\sigma,\eps\in (0,1]$, there exist constants~$\kappa = \kappa(d, \sigma)>0$ and~$R = R(d,\sigma,\eps)\in\N$ such that the following holds.
For every natural rank function~$\rank: \Pol_{\leq d}(\F^{\infty}, \F^k) \to \R_+$ and every map~$\phi \in \Pol_{\leq d}(\F^n, \F^k)$ with~$\rank(\phi) \geq R$, we have
\beqn
\Pr_{I\sim [n]_{\sigma}}\big[ \rank(\phi_{|I}) \geq \kappa\cdot \rank(\phi)\big] \geq 1 - \eps.
\eeqn
\end{theorem}
\setcounter{subsection}{2}
\subsection{Notation}
\label{sec:prelims}

Here we collect some notation that will be used throughout the paper.
Given sets  $I_1, \dots, I_d$, we write $I_{[d]} = I_1\times \cdots\times I_d$.
For a set~$J$ and some~$\sigma \in [0,1]$, we denote by~$\pi_{\sigma}^J$ the product distribution on~$\{0,1\}^J$ where each coordinate is independently set to $1$ with probability~$\sigma$, or to $0$ with probability~$1-\sigma$;
when~$J=[n]$, we write simply~$\pi_{\sigma}^n$.
We write $J_\sigma$ for the probability distribution over subsets of~$J$ where each element is present independently with probability~$\sigma$.
To denote that a random variable~$X$ is distributed according to a distribution~$\mu$, we write $X\sim \mu$.
\section{Tensors}
\label{sec:tensors}

Recall that, for a field $\F$ and integer $d\geq 2$, we denote by $(\F^\infty)^{\otimes d}$ the set of all $d$-tensors of finite support over $\F$.
All our results (including the values of the implied constants) will hold independently of the field considered, so we will always denote it by $\F$ without further comment.

The notions of tensor rank we will consider here are those called \emph{natural rank functions} as defined below:

\begin{definition}[Natural rank]\label{def:nat_rank}
We say that $\rank: (\F^{\infty})^{\otimes d} \to \R_+$ is a natural rank function if $\rank({\bf 0}) = 0$ and it satisfies the following properties:
\begin{enumerate}
    \item Sub-additivity: \\
    $\rank(T+S) \leq \rank(T) + \rank(S)$ for all~$T, S\in (\F^{\infty})^{\otimes d}$.
    \item Monotonicity under restrictions: \\
    $\rank\big(T_{|I_{[d]}}\big) \leq \rank(T)$ for all~$T\in (\F^{\infty})^{\otimes d}$ and all sets~$I_1, \dots, I_d \subset \N$.
    \item Restriction Lipschitz property: \\
    $\rank\big(T_{|J_{[d]}}\big) \leq \rank\big(T_{|I_{[d]}}\big) + \sum_{i=1}^d |J_i \setminus I_i|$
    for all~$T\in (\F^{\infty})^{\otimes d}$ and all sets~$I_1\subseteq J_1, \dots, I_d\subseteq J_d$.
    \end{enumerate}
\end{definition}
\setcounter{theorem}{2}
\subsection{Concentration for monotone sub-additive functions}

The main step in our proof of Theorem~\ref{thm:random_tensor} is a type of concentration inequality for monotone sub-additive functions on the hypercube.

In order to obtain such a result we will consider a notion of \emph{two-point distance} from sets $A\subseteq \bset{n}$, which intuitively measures how many coordinates of some given point $x$ cannot be captured by any two elements of~$A$.

\begin{definition}
Given a point $x\in \bset{n}$ and a set $A \subseteq \bset{n}$, the two-point distance between $x$ and $A$ is
\beqn
h_2(x; A) = \min_{y, z \in A} \big|\big\{i \in [n] \st x_i \neq y_i \And x_i \neq z_i\big\}\big|.
\eeqn
\end{definition}

Note that this ``distance'' $h_2(x; A)$ can be zero even if $x \notin A$;
for instance, $h_2\big(00;\, \{01, 10\}\big) = 0$.
The reason for considering this notion is that one gets much better concentration for $h_2$ than one does for the usual Hamming distance.
This is shown by the following result, which is a special case of an inequality of Talagrand~\cite[Theorem 3.1.1]{Talagrand:1995}:

\begin{theorem}\label{thm:Talagrand}
For any parameter $\sigma\in [0,1]$ and set $A \subseteq \bset{n}$, we have that
\beqn
\Pr_{x\sim \pi_\sigma^n}\big[h_2(x; A) \geq k\big] \leq 2^{-k} \pi_\sigma^n(A)^{-2} \quad \text{for all $k\geq 0$.}
\eeqn
\end{theorem}

The next lemma is the key result needed for proving our main theorem for tensors, and it deals more abstractly with monotone, sub-additive Lipschitz functions on the hypercube.
We endow~$\bset{n}$ with the usual partial order, where $x\leq y$ if the support of~$x$ is contained in the support of~$y$.
A function $f:\bset{n} \to \R$ is monotone if $f(x) \leq f(y)$ whenever $x\leq y$, and it is sub-additive if $f(x+y) \leq f(x) + f(y)$ for all $x,y\in\bset{n}$ with  disjoint supports.
Finally,~$f$ is $1$-Lipschitz if
\beqn
\big|f(x) - f(y)\big| \leq \big|\big\{i \in [n] \st x_i \neq y_i\big\}\big| \quad \text{for all $x, y \in \bset{n}$}.
\eeqn
For a string $x\in \bset{n}$ and set $I\subseteq[n]$, we let~$x_I\in \bset{n}$ be the string that equals~$x$ on the indices in~$I$ and is zero elsewhere.
The lemma that follows and its proof are inspired by an argument of Schechtman~\cite[Corollary~12]{Schechtman2003}.

\begin{lemma}[Concentration inequality] \label{lem:subadditive}
Let $f:\bset{n} \to \R_+$ be a monotone, sub-additive $1$-Lipschitz function such that $f({\bf 0}) =0$ and $f({\bf 1}) = r$.
Then,
\begin{align*}
    \Pr_{x\sim \pi_{\sigma}^n} \big[f(x) \leq \sigma r/4\big] &\leq \sqrt{\frac{3}{2\sigma}}\, 2^{-\sigma r/12} \quad &\text{if $\,0 < \sigma < 1/2$,} \\
    \Pr_{x\sim \pi_{\sigma}^n} \big[f(x) \leq r/6\big] &\leq \sqrt{2}\, 2^{-r/12} &\text{if $\,1/2 \leq \sigma \leq 1$.}
\end{align*}
\end{lemma}

\begin{proof}
We first prove the first inequality above.
Let $\sigma\in (0, 1/2)$, and consider the set
\beqn
A = \big\{ x\in \bset{n} \st f(x)\leq \sigma r/4 \big\}.
\eeqn
By Talagrand's Inequality (Theorem~\ref{thm:Talagrand}) we have
\beqn
\Pr_{x\sim \pi_{\sigma}^n}\big[ h_2(x; A)\geq \sigma r/6\big] \leq 2^{-\sigma r/6} \,\Pr_{x\sim \pi_{\sigma}^n}\big[ f(x)\leq \sigma r/4\big]^{-2},
\eeqn
so to finish the proof it suffices to show that
\beqn
\Pr_{x\sim \pi_{\sigma}^n}\big[ h_2(x; A)\geq \sigma r/6\big] \geq 2\sigma/3.
\eeqn

We claim that
\beq \label{eq:h2inclusion}
\big\{ x\in \bset{n} \st f(x)\geq 2\sigma r/3 \big\} \subseteq \big\{ x\in \bset{n} \st h_2(x; A)\geq \sigma r/6 \big\}.
\eeq
Indeed, if $h_2(x; A) < \sigma r/6$ then there are $y, z\in A$ such that
\beqn
\big|\big\{i \in [n] \st x_i \neq y_i \And x_i \neq z_i\big\}\big| < \sigma r/6.
\eeqn
Denote $I = \big\{i \in [n] \st x_i \neq y_i \And x_i \neq z_i\big\}$, $J = \big\{i \in [n] \st x_i = y_i\big\}$ and $J' = \big\{i \in [n] \st x_i \neq y_i \And x_i = z_i\big\}$;
note that these sets partition~$[n]$, and by assumption $|I| < \sigma r/6$.
We then have
\begin{align*}
    f(x) &\leq f(x_I + x_J) + f(x_{J'}) \\
    &= f(x_I + y_J) + f(z_{J'}) \\
    &\leq |I| + f(y_J) + f(z_{J'}) \\
    &\leq |I| + f(y) + f(z) \\
    &< 2\sigma r/3,
\end{align*}
so $\big\{h_2(x; A) < \sigma r/6\big\} \subseteq \big\{f(x) < 2\sigma r/3\big\}$ and inclusion \eqref{eq:h2inclusion} follows.
It thus suffices to show that $\Pr_{x\sim \pi_{\sigma}^n}\big[ f(x)\geq 2\sigma r/3\big] \geq 2\sigma/3$.

We will next prove that, for any integer $k\geq 1$, we have
\beq \label{eq:coupling}
\Pr_{x\sim \pi_{1/k}^n}\big[ f(x)\geq r/k\big] \geq 1/k.
\eeq
This is done by a simple \emph{coupling argument}, which will be important for us again later on.
Consider a uniformly random ordered $k$-partition $(I_1, \dots, I_k)$ of $[n]$;
thus the $k$ sets $I_i$ are pairwise disjoint and have union $[n]$, with each of the $k^n$ possible such $k$-tuples having the same probability.
Denoting by $\1_I$ the indicator function of set $I$, we have
\beqn
r = f(\1_{[n]}) = f(\1_{I_1} + \dots + \1_{I_k}) \leq \sum_{i=1}^k f(\1_{I_i}),
\eeqn
so $f(\1_{I_i}) \geq r/k$ holds for at least one $i\in [k]$ in \emph{every} ordered partition $(I_1, \dots, I_k)$.
By symmetry, it follows that $\Pr \big[f(\1_{I_1}) \geq r/k \big] \geq 1/k$.
Since the marginal distribution of $\1_{I_1}$ (and every other $\1_{I_i}$) is precisely $\pi_{1/k}^n$, we conclude that
\begin{equation*}
\Pr_{x\sim \pi_{1/k}^n} \big[f(x) \geq r/k\big] = \Pr \big[f(\1_{I_1}) \geq r/k \big] \geq 1/k,
\end{equation*}
as wished.

Now let $k = \lceil 1/\sigma \rceil$.
Since $0 < \sigma < 1/2$, we have $2\sigma/3 \leq 1/k \leq \sigma$, and so by monotonicity of $f$
\begin{align*}
    \Pr_{x\sim \pi_{\sigma}^n}\big[ f(x)\geq 2\sigma r/3\big] &\geq \Pr_{x\sim \pi_{1/k}^n}\big[ f(x)\geq 2\sigma r/3\big] \\
    &\geq \Pr_{x\sim \pi_{1/k}^n}\big[ f(x)\geq r/k\big].
\end{align*}
From inequality~\eqref{eq:coupling} we then conclude that
\beqn
\Pr_{x\sim \pi_{\sigma}^n}\big[ f(x)\geq 2\sigma r/3\big] \geq 1/k \geq 2\sigma/3,
\eeqn
finishing the proof of the first inequality in the statement of the lemma.

The second inequality is proven using the same arguments, now applied to the parameter $\sigma = 1/2$ and the set $A = \big\{ x\in \bset{n} \st f(x)\leq r/6 \big\}$.
By the same argument as above,
\begin{equation*}
    \big\{x\in \bset{n}\st f(x) \geq r/2\big\} \subseteq \big\{x\in \bset{n}\st h_2(x;A) \geq r/6\big\}.
\end{equation*}
Talagrand's inequality and coupling imply that
\begin{align*}
    2^{-r/6}\Pr_{x\sim\pi_{1/2}^n}\big[f(x) \leq r/6\big]^{-2} &\geq \Pr_{x\sim \pi_{1/2}^n}\big[h_2(x;A) \geq r/6\big]\\
    &\geq \Pr_{x\sim\pi_{1/2}^n}\big[f(x) \geq r/2\big]\\
    &\geq \frac{1}{2}.
\end{align*}
Rearranging then gives the result for~$\sigma =  1/2$;
we obtain slightly better bounds since in this case $1/\sigma = 2$ is an integer, and so no rounding errors occur.
By monotonicity of $f$, the same bound continues to hold for all $\sigma > 1/2$.
\end{proof}
\setcounter{section}{2}
\section{Polynomial maps}
\label{sec:polynomials}

Next we consider the setting of polynomials and higher-dimensional polynomial maps.
Recall that $\Pol_{\leq d}(\F^n, \F^k)$ is the space of all polynomial maps~$\phi: \F^n \to \F^k$ of degree at most~$d$, and $\Pol_{\leq d}(\F^{\infty}, \F^k) = \bigcup_{n\in \N} \Pol_{\leq d}(\F^n, \F^k)$ is the space of all $k$-dimensional polynomial maps over~$\F$ of degree at most~$d$.

\begin{definition}[Natural rank]
We say that $\rank: \Pol_{\leq d}(\F^{\infty}, \F^k) \to \R_+$ is a natural rank function if $\rank({\bf 0}) = 0$ and it satisfies the following properties:
\begin{enumerate}
    \item Symmetry: \\
    $\rank(\phi) = \rank(-\phi)$ for all~$\phi\in \Pol_{\leq d}(\F^{\infty}, \F^k)$.
    \item Sub-additivity: \\
    $\rank(\phi + \gamma) \leq \rank(\phi) + \rank(\gamma)$ for all~$\phi, \gamma\in \Pol_{\leq d}(\F^{\infty}, \F^k)$.
    \item Monotonicity under restrictions: \\
    $\rank(\phi_{|I}) \leq \rank(\phi)$ for all~$\phi\in \Pol_{\leq d}(\F^{\infty}, \F^k)$ and all sets~$I \subset \N$.
    \item Restriction Lipschitz property: \\
    $\rank(\phi_{|I \cup J}) \leq \rank(\phi_{|I}) + |J|$ for all~$\phi\in \Pol_{\leq d}(\F^{\infty}, \F^k)$ and all sets~$I, J \subset \N$.
\end{enumerate}
\end{definition}
\setcounter{equation}{3}
\begin{theorem}[Theorem~\ref{thm:random_poly} restated]
For every~$d\in \N$ and~$\sigma,\eps\in (0,1]$, there exist constants~$\kappa = \kappa(d, \sigma)>0$ and~$R = R(d,\sigma,\eps)\in\N$ such that the following holds.
For every natural rank function~$\rank: \Pol_{\leq d}(\F^{\infty}, \F^k) \to \R$ and every map~$\phi \in \Pol_{\leq d}(\F^n, \F^k)$ with~$\rank(\phi) \geq R$, we have
\beqn
\Pr_{I\sim [n]_{\sigma}}\big[ \rank(\phi_{|I}) \geq \kappa\cdot \rank(\phi)\big] \geq 1 - \eps.
\eeqn
\end{theorem}

The proof of this theorem proceeds differently from the proof of Theorem~\ref{thm:random_tensor}, which gives the analogous result for tensors.
The reason for this is that the natural choice of function $f:\bset{n}\to \R_+$ to which one might apply the concentration inequality in Lemma~\ref{lem:subadditive}, given by $f(\one_I) = \rank(\phi_{|I})$ (for some rank function~$\rank$), might fail to be sub-additive as a function on the boolean hypercube while~$\rank$ is sub-additive as a natural rank function.
For instance, consider (in dimension~$k=1$) a $2n$-variate polynomial $p(x_1,\dots,x_n,y_1,\dots,y_n)$ such that each monomial contains both an~$x$ variable and a~$y$ variable.
Then the rank of~$p$ restricted only to its~$x$ variables is zero, and similarly for the restriction to the~$y$ variables, while~$\rank(p)$ can be of order~$\Theta(n)$.
This shows that~$\rank(p_{|I \cup J})$ can be much larger than $\rank(p_{|I}) + \rank(p_{|J})$, and the argument used in the tensor case breaks down.


\subsection{The proof in expectation}

The first step in our proof of Theorem~\ref{thm:random_poly} is to show that the desired result is true \emph{in expectation}, rather than with high probability.
More precisely, we wish to show an inequality of the form
\beqn
\Exp_{I\sim [n]_{\sigma}} \rank(\phi_{|I}) \geq c(\sigma, d) \rank(\phi)
\eeqn
which is valid for all degree-$d$ maps $\phi\in \Pol_{\leq d}(\F^n, \F^k)$ and all natural rank functions.

Towards proving the inequality above by induction on the degree, it will be convenient to generalize it to polynomial maps over \emph{commutative rings} rather than only fields.
The reason for doing so is that, by considering some subset of variables $\{x_a: a\in A\}$ to be constants, we might be able to decrease the degree of the associated polynomial map over the remaining variables $\{x_b: b\notin A\}$.

More formally, given a polynomial map $\phi\in \Pol_{\leq d}(\F^n, \F^k)$ and two disjoint subsets $A, B \subset [n]$, denote by $\phi[A, B]$ the sum of the monomials in~$\phi$ involving at least one variable from each~$A$ and~$B$.
Note that $\phi[A, B]$ can be regarded as a polynomial map in three different ways:
either as having coefficients in~$\F$ and variables in $\{x_i: i\in [n]\}$;
or having coefficients in the ring $\F[x_b: b\in B]$ and variables in $\{x_a: a\in A\}$;
or having coefficients in $\F[x_a: a\in A]$ and variables in $\{x_b: b\in B\}$.
The advantage of these last two representations is that the degree of $\phi[A, B]$ \emph{decreases}, since every monomial will contain formal variables which are now elements of the ring (and are thus regarded as constants).

For a commutative ring~$R$, denote by $R[x_1, \dots, x_n]_{\leq d}$ the set of formal polynomials in variables $x_1, \dots, x_n$ with coefficients over~$R$ and degree at most~$d$.
For a bipartition $[n] = A\uplus B$, define the polynomial ring $R_A = R[x_a\st a\in A]$ and consider the natural identification of $R[x_1,\dots,x_n]$ and $R_A[x_b\st b\in B]$.
Under this identification, a function $\rank: (R[x_1, \dots, x_n])^k \to \R_+$ induces a function $\rank:(R_A[x_b\st b\in B])^k\to \R_+$ in a way that preserves symmetry, sub-additivity and monotonicity under restrictions.


\begin{lemma}[Linear expectation] \label{lem:expectation}
For every~$\sigma\in (0, 1]$ and every integer~$d \geq 1$ there exists a constant~$c(\sigma, d) > 0$ such that the following holds.
Let~$n,k$ be positive integers,~$R$ be a commutative ring and $\rank: (R[x_1, \dots, x_n]_{\leq d})^k \to \R_+$ be symmetric, sub-additive and monotone under restrictions.
Then, for any map $\phi\in (R[x_1, \dots, x_n]_{\leq d})^k$ we have
\beqn
\Exp_{J\sim [n]_{\sigma}} \rank(\phi_{|J}) \geq c(\sigma, d) \rank(\phi).
\eeqn
\end{lemma}

\begin{proof}
For simplicity, we will prove the result in the special case where $\sigma=1/2$.
Since $\phi_{|I\cap J} = (\phi_{|I})_{|J}$ and $I\cap J \sim [n]_{1/2^{j+1}}$ when $I\sim [n]_{1/2}$, $J\sim [n]_{1/2^j}$, the general case easily follows from this special case by taking constant
\beqn
c(\sigma, d) = c(1/2, d)^{\lceil\log(1/\sigma)\rceil}
\eeqn
and using monotonicity under restrictions.

Denote $c(d) := c(1/2, d)$ for ease of notation.
The proof will proceed by induction on the degree~$d$, with the base case where $d=0$ holding trivially with $c(0) = 1$.
For the inductive step, let
\beqn
c(d) = \frac{c(d-1)^2}{20}
\eeqn
and consider two different possibilities.

First, suppose that for every bipartition $[n] = A\uplus B$ we have that
\beqn
\max \big\{\rank(\phi_{|A}),\, \rank(\phi_{|B})\big\} \geq 2c(d) \rank(\phi).
\eeqn
Using a similar coupling argument as was used in the proof of Lemma~\ref{lem:subadditive}, we get that
\begin{align*} \Pr_{I\sim[n]_{1/2}}\big[\rank\big(\phi_{|I}\big) \geq 2c(d)\rank(\phi)\big] \geq \frac{1}{2},
\end{align*}
which implies
\begin{equation*}
    \Exp_{I\sim [n]_{1/2}}\rank\big(\phi_{|I}\big) \geq c(d)\rank(\phi)
\end{equation*}
as desired.

The second possibility is that there exists a bipartition $[n] = A\uplus B$ such that
\beq \label{eq:second_case}
\rank(\phi_{|A}) < 2c(d) \rank(\phi) \quad \text{and} \quad \rank(\phi_{|B}) < 2c(d) \rank(\phi).
\eeq
Fix such a bipartition for the rest of the proof.
For any given subsets $I\subseteq A$, $J\subseteq B$, let $\phi[I, J]$ be the sum of the monomials in~$\phi$ involving at least one variable from each~$I$ and~$J$.
It is easy to see that
\beqn
\phi[I, J] = \phi_{|I\cup J} - \phi_{|I} - \phi_{|J} + \phi(0).
\eeqn
Define the commutative rings $R_A := R[x_a: a\in A]$ and $R_B := R[x_b: b\in B]$, and note that $\phi[I, J]$ can be seen both as an element of $(R_A[x_b: b\in B])^k$ and as an element of $(R_B[x_a: a\in A])^k$.
Moreover, the degree of $\phi[I, J]$ is at most $d-1$ in each of these representations.

From the identity $\phi = \phi_{|A} + \phi_{|B} - \phi(0) + \phi[A, B]$ and sub-additivity, it follows that
\beqn
\rank(\phi[A, B]) \geq \rank(\phi) - \rank(\phi_{|A}) - \rank(\phi_{|B}) - \rank(-\phi(0)).
\eeqn
Since $\rank(-\phi(0)) = \rank(\phi(0)) = \rank(\phi_{|\emptyset})$, we conclude from monotonicity under restrictions and our assumption~\eqref{eq:second_case} that
\beq\label{eq:phiAB}
\rank(\phi[A, B]) > \rank(\phi) - 6 c(d) \rank(\phi) > \rank(\phi)/2.
\eeq
In an analogous way, we obtain the inequalities
\begin{align} \label{eq:phi_IJ}
    \rank(\phi_{|I\cup J}) &\geq \rank(\phi[I, J]) - \rank(\phi_{|I}) - \rank(\phi_{|J}) - \rank(\phi(0)) \\
    &> \rank(\phi[I, J]) - 6 c(d) \rank(\phi) \notag
\end{align}
valid for all $I\subseteq A$, $J\subseteq B$.

Applying the inductive hypothesis to $\phi[A, B] \in (R_B[x_a: a\in A]_{\leq d-1})^k$ we obtain
\beqn
\Exp_{I\sim A_{1/2}} \rank(\phi[I, B]) = \Exp_{I\sim A_{1/2}} \rank(\phi[A, B]_{|I}) \geq c(d-1) \rank(\phi[A, B]).
\eeqn
Moreover, for any fixed $I\subseteq A$, the inductive hypothesis applied to the map $\phi[I, B] \in (R_A[x_b: b\in B]_{\leq d-1})^k$ gives
\beqn
\Exp_{J\sim B_{1/2}} \rank(\phi[I, J]) = \Exp_{J\sim B_{1/2}} \rank(\phi[I, B]_{|J}) \geq c(d-1) \rank(\phi[I, B]).
\eeqn
Combining the two inequalities above we conclude that
\beqn
\Exp_{I\sim A_{1/2},\, J\sim B_{1/2}} \rank(\phi[I, J]) \geq c(d-1)^2 \rank(\phi[A, B]),
\eeqn
and thus by~\eqref{eq:phiAB} we have
$\Exp_{I\sim A_{1/2},\, J\sim B_{1/2}} \rank(\phi[I, J]) > c(d-1)^2 \rank(\phi)/2$.
Together with equation~$\eqref{eq:phi_IJ}$, we conclude that
\begin{align*}
    \Exp_{I\sim A_{1/2},\, J\sim B_{1/2}}\rank(\phi_{|I\cup J})
    &> \Exp_{I\sim A_{1/2},\, J\sim B_{1/2}} \rank(\phi[I, J]) - 6c(d) \rank(\phi) \\
    &> c(d-1)^2 \rank(\phi)/2 - 6c(d) \rank(\phi) \\
    &> c(d) \rank(\phi),
\end{align*}
which is precisely what we wanted to prove.
\end{proof}




\subsection{Monotone functions on the hypercube and boosting}

Our next result is a lemma which allows us to boost the probability of some events (such as having high rank under random restrictions) from~$\eps$ to~$1-\eps$ by paying a relatively small price.

It will again be convenient to take a more abstract approach and deal with Boolean functions rather than restrictions of polynomials.
For a Boolean function~$g: \{0,1\}^n \to \{0,1\}$, the total influence of~$g$ under distribution~$\pi_{\sigma}^n$ is given by
\beqn
\mathbf{I}^{(\sigma)}(g)
=
\sum_{i=1}^n \Exp_{x\sim\pi_\sigma^n}\big|g(x) - g(x^i)\big|,
\eeqn
where~$x^i$ differs from~$x$ only in the~$i$th coordinate.
Denote by~$\mu_{\sigma}[g] := \Exp_{x\sim \pi_{\sigma}^n} \big[g(x)\big]$ its expectation.

\begin{lemma}[Boosting lemma] \label{lem:boost}
For every~$\sigma>0$ there is a constant~$C_{\sigma}>0$ such that the following holds.
Let~$n\in \N$, $f: \{0,1\}^n \to \R_+$ be a monotone~$1$-Lipschitz function and~$\eps\in (0, 1/2]$.
If~$r\geq 1$ satisfies
\beqn
\Pr_{x\sim \pi_{\sigma}^n} \big[f(x) \geq r\big] > \eps,
\eeqn
then~$\Pr_{x\sim \pi_{1.01\sigma}^n} \big[f(x) \geq r - C_{\sigma}^{1/\eps^2}\big] > 1-\eps$.
\end{lemma}

\begin{proof}
Consider the Boolean function~$g(x) := 1 \big[f(x) \geq r\big]$.
Note that there exists~$q \in [\sigma,\, 1.01\sigma]$ such that~$\mathbf{I}^{(q)}(g) \leq 100/\sigma$, as otherwise by the Margulis-Russo formula we would have
\beqn
\mu_{1.01\sigma}[g] = \mu_{\sigma}[g] + \int_{\sigma}^{1.01\sigma} \frac{d\mu_q[g]}{dq} dq \geq \int_{\sigma}^{1.01\sigma} \mathbf{I}^{(q)}(g) dq > 1.
\eeqn

Let~$\gamma \in (0, \eps/2)$ be a constant to be chosen later.
By Friedgut's Junta Theorem~\cite{Friedgut1998}, there is a~$C(q)^{100/\gamma \sigma}$-junta~$h: \{0,1\}^n \to \{0,1\}$ such that
\beqn
\Pr_{x\sim \pi_q^n} \big[g(x) \neq h(x)\big] < \gamma.
\eeqn
Let~$J \subseteq [n]$, $|J| \leq C(q)^{100/\gamma \sigma}$, be the set of variables on which~$h$ depends.
Then
\begin{align*}
    \mu_q[h] \geq \mu_q[g] - \gamma \geq \mu_{\sigma}[g] - \gamma > \eps/2,
\end{align*}
which implies~$\Pr_{z\sim \pi_q^J} \big[h(z) = 1\big] > \eps/2$.
Since
\beqn
\Exp_{z\sim \pi_q^J} \Pr_{y\sim \pi_q^{J^c}} \big[g(y,z) \neq h(z)\big] < \gamma,
\eeqn
it follows that
\beqn
\Pr_{z\sim \pi_q^J} \big[\Pr_{y\sim \pi_q^{J^c}} \big[g(y,z) \neq h(z)\big] \geq \eps\big] < \gamma/\eps.
\eeqn
Taking~$\gamma = \eps^2/2$, we conclude there exists~$z\in \{0,1\}^J$ such that~$h(z)=1$ and
\beqn
    \Pr_{y\sim \pi_q^{J^c}} \big[g(y,z) = 0\big] = \Pr_{y\sim \pi_q^{J^c}} \big[g(y,z) \neq h(z)\big] < \eps.
\eeqn

Since~$f$ is~$1$-Lipschitz by assumption, we have
\beqn
    f(y,z) \geq r \implies f(y,z') \geq r - |J| \quad \forall z'\in \{0,1\}^J,
\eeqn
and thus by monotonicity
\begin{align*}
    \Pr_{x\sim \pi_{1.01\sigma}^n} \big[f(x) \geq r - |J|\big]
    &\geq \Pr_{x\sim \pi_q^n} \big[f(x) \geq r - |J|\big] \\
    &= \Exp_{z'\sim \pi_q^J} \Pr_{y\sim \pi_q^{J^c}} \big[f(y,z') \geq r - |J|\big] \\
    &\geq \max_{z\in \{0,1\}^J} \Pr_{y\sim \pi_q^{J^c}} \big[f(y,z) \geq r\big] \\
    &> 1-\eps.
\end{align*}
This is precisely what we wanted to prove, with constant
$$C_{\sigma} = \sup_{q\in [\sigma,\, 1.01\sigma]} C(q)^{200/\sigma}.$$
Note that the above supremum is finite:
in \cite[Section~10.3]{ODonnell:2014} it is shown that we can take
\beqn
C(q) = \max\big\{ q^{-c},\, (1-q)^{-c}\big\}
\eeqn
for some absolute constant $c>0$.
\end{proof}


\subsection{The Random Restriction Theorem}

Our main result, Theorem~\ref{thm:random_poly}, follows easily from the lemmas given above.
Recall that we wish to show that
\beqn
\Pr_{J\sim [n]_{\sigma}}\big[ \rank(\phi_{|J}) \geq \kappa(d, \sigma)\cdot \rank(\phi)\big] \geq 1 - \eps
\eeqn
whenever~$\rank(\phi) \geq R(d, \sigma, \eps)$, for some well-chosen constants~$R(d, \sigma, \eps)$ and $\kappa(d, \sigma) > 0$.

\begin{proof}[ of Theorem~\ref{thm:random_poly}]
Applying Lemma \ref{lem:expectation} with~$\sigma$ substituted by~$0.9\sigma$, we obtain
\beqn
\Exp_{J\sim [n]_{0.9\sigma}} \rank(\phi_{|J}) \geq c(0.9\sigma, d) \rank(\phi).
\eeqn
Since~$\rank(\phi_{|J}) \leq \rank(\phi)$ for all subsets~$J \subseteq [n]$, we conclude that
\beqn
\Pr_{J\sim [n]_{0.9\sigma}} \bigg[ \rank(\phi_{|J}) \geq \frac{c(0.9\sigma, d)}{2} \rank(\phi) \bigg] \geq \frac{c(0.9\sigma, d)}{2}.
\eeqn

By possibly decreasing~$\eps$ a little, we may assume that~$\eps < c(0.9\sigma, d)/2$.
Denote
\beqn
\kappa(d, \sigma) = \frac{c(0.9\sigma, d)}{4} \quad \text{and} \quad R(d, \sigma, \eps) = \frac{4 C_{0.9\sigma}^{1/\eps^2}}{c(0.9\sigma, d)},
\eeqn
where~$C_{0.9\sigma} > 0$ is the constant guaranteed by the boosting lemma, Lemma \ref{lem:boost}.
Applying the boosting lemma to the function~$J \mapsto \rank(\phi_{|J})$ with~$r = c(0.9\sigma, d) \rank(\phi)/2$, we get
\begin{align*}
    \Pr_{J\sim [n]_{0.9\sigma}} \big[ \rank(\phi_{|J}) \geq r \big] &\geq c(0.9\sigma, d)/2 > \eps \\
    &\implies \Pr_{J\sim [n]_{\sigma}} \big[ \rank(\phi_{|J}) \geq r - C_{0.9\sigma}^{1/\eps^2} \big] > 1-\eps.
\end{align*}
Since~$c(0.9\sigma, d) \rank(\phi)/2 - C_{0.9\sigma}^{1/\eps^2} \geq \kappa(\sigma, d) \rank(\phi)$ whenever $\rank(\phi) \geq R(d, \sigma, \eps)$, the result follows.
\end{proof}
\begin{thebibliography}{10}

\bibitem{AnanyanH:2020}
Tigran Ananyan and Melvin Hochster.
\newblock Small subalgebras of polynomial rings and {S}tillman's conjecture.
\newblock {\em J. Amer. Math. Soc.}, 33:291--309, 2020.
\newblock \href {https://doi.org/10.1090/jams/932} {\path{doi:10.1090/jams/932}}.

\bibitem{BCCGNSU}
Jonah Blasiak, Thomas Church, Henry Cohn, Joshua~A. Grochow, Eric Naslund, William~F. Sawin, and Chris Umans.
\newblock On cap sets and the group-theoretic approach to matrix multiplication.
\newblock {\em Discrete Anal.}, pages Paper No. 3, 27, 2017.
\newblock \href {https://doi.org/10.19086/da.1245} {\path{doi:10.19086/da.1245}}.

\bibitem{BlatterDR:2022}
Andreas Blatter, Jan Draisma, and Filip Rupniewski.
\newblock A tensor restriction theorem over finite fields, 2022.
\newblock arXiv:2211.12319.
\newblock \href {https://doi.org/10.48550/ARXIV.2211.12319} {\path{doi:10.48550/ARXIV.2211.12319}}.

\bibitem{Briet:2021}
Jop Bri\"{e}t.
\newblock Subspaces of tensors with high analytic rank.
\newblock {\em Online J. Anal. Comb.}, (16):Paper No. 6, 9, 2021.

\bibitem{BrietBCSN}
Jop Briët, Harry Buhrman, Davi Castro-Silva, and Niels M.~P. Neumann.
\newblock Noisy decoding by shallow circuits with parities: classical and quantum, 2023.
\newblock \href {https://doi.org/10.48550/ARXIV.2302.02870} {\path{doi:10.48550/ARXIV.2302.02870}}.

\bibitem{CohenM21}
Alex Cohen and Guy Moshkovitz.
\newblock Structure vs.\ randomness for bilinear maps.
\newblock In {\em S{TOC} '21---{P}roceedings of the 53rd {A}nnual {ACM} {SIGACT} {S}ymposium on {T}heory of {C}omputing}, pages 800--808. ACM, New York, 2021.
\newblock \href {https://doi.org/10.1145/3406325.3451007} {\path{doi:10.1145/3406325.3451007}}.

\bibitem{CohenM21b}
Alex Cohen and Guy Moshkovitz.
\newblock Partition and analytic rank are equivalent over large fields.
\newblock {\em Duke Math. J.}, 172(12):2433--2470, 2023.
\newblock \href {https://doi.org/10.1215/00127094-2022-0086} {\path{doi:10.1215/00127094-2022-0086}}.

\bibitem{CrootLP2017}
Ernie Croot, Vsevolod~F. Lev, and P\'{e}ter~P\'{a}l Pach.
\newblock Progression-free sets in~{$\Bbb Z^n_4$} are exponentially small.
\newblock {\em Ann. of Math. (2)}, 185(1):331--337, 2017.
\newblock \href {https://doi.org/10.4007/annals.2017.185.1.7} {\path{doi:10.4007/annals.2017.185.1.7}}.

\bibitem{DerksenES:2017}
Harm Derksen, Rob~H. Eggermont, and Andrew Snowden.
\newblock Topological noetherianity for cubic polynomials.
\newblock {\em Algebra Number Theory}, 11(9):2197--2212, 2017.
\newblock \href {https://doi.org/10.2140/ant.2017.11.2197} {\path{doi:10.2140/ant.2017.11.2197}}.

\bibitem{Draisma:2019}
Jan Draisma.
\newblock Topological {N}oetherianity of polynomial functors.
\newblock {\em J. Amer. Math. Soc.}, 32(3):691--707, 2019.
\newblock \href {https://doi.org/10.1090/jams/923} {\path{doi:10.1090/jams/923}}.

\bibitem{EllenbergG2017}
Jordan~S. Ellenberg and Dion Gijswijt.
\newblock On large subsets of {$\Bbb F^n_q$} with no three-term arithmetic progression.
\newblock {\em Ann. of Math. (2)}, 185(1):339--343, 2017.
\newblock \href {https://doi.org/10.4007/annals.2017.185.1.8} {\path{doi:10.4007/annals.2017.185.1.8}}.

\bibitem{FoxL17}
Jacob Fox and L\'{a}szl\'{o}~Mikl\'{o}s Lov\'{a}sz.
\newblock A tight bound for {G}reen's arithmetic triangle removal lemma in vector spaces.
\newblock {\em Adv. Math.}, 321:287--297, 2017.
\newblock \href {https://doi.org/10.1016/j.aim.2017.09.037} {\path{doi:10.1016/j.aim.2017.09.037}}.

\bibitem{Friedgut1998}
Ehud Friedgut.
\newblock Boolean functions with low average sensitivity depend on few coordinates.
\newblock {\em Combinatorica}, 18(1):27--35, 1998.
\newblock \href {https://doi.org/10.1007/PL00009809} {\path{doi:10.1007/PL00009809}}.

\bibitem{GowersK2022}
W.~T. Gowers and Thomas Karam.
\newblock Equidistribution of high-rank polynomials with variables restricted to subsets of $\mathbb{F}_p$, 2022.
\newblock arXiv:2209.04932.
\newblock \href {https://doi.org/10.48550/ARXIV.2209.04932} {\path{doi:10.48550/ARXIV.2209.04932}}.

\bibitem{GowersW2011}
W.~T. Gowers and J.~Wolf.
\newblock Linear forms and higher-degree uniformity for functions on {$\Bbb F^n_p$}.
\newblock {\em Geom. Funct. Anal.}, 21(1):36--69, 2011.
\newblock \href {https://doi.org/10.1007/s00039-010-0106-3} {\path{doi:10.1007/s00039-010-0106-3}}.

\bibitem{GreenT2009}
Ben Green and Terence Tao.
\newblock The distribution of polynomials over finite fields, with applications to the {G}owers norms.
\newblock {\em Contrib. Discrete Math.}, 4(2):1--36, 2009.

\bibitem{HagerupRub:1990}
Torben Hagerup and Christine R{\"u}b.
\newblock A guided tour of {Chernoff} bounds.
\newblock {\em Information Processing Letters}, 33(6):305--308, 1990.
\newblock \href {https://doi.org/10.1016/0020-0190(90)90214-I} {\path{doi:10.1016/0020-0190(90)90214-I}}.

\bibitem{HaramatyS2010}
Elad Haramaty and Amir Shpilka.
\newblock On the structure of cubic and quartic polynomials.
\newblock In {\em S{TOC}'10---{P}roceedings of the 2010 {ACM} {I}nternational {S}ymposium on {T}heory of {C}omputing}, pages 331--340. ACM, New York, 2010.
\newblock \href {https://doi.org/10.1145/1806689.1806736} {\path{doi:10.1145/1806689.1806736}}.

\bibitem{Hatami2016}
Pooya Hatami.
\newblock On the structure of quintic polynomials.
\newblock In {\em Approximation, randomization, and combinatorial optimization. {A}lgorithms and techniques}, volume~60 of {\em LIPIcs. Leibniz Int. Proc. Inform.}, pages Art. No. 33, 18. Schloss Dagstuhl. Leibniz-Zent. Inform., Wadern, 2016.
\newblock \href {https://doi.org/10.4230/LIPIcs.APPROX-RANDOM.2016.33} {\path{doi:10.4230/LIPIcs.APPROX-RANDOM.2016.33}}.

\bibitem{Karam2022}
Thomas Karam.
\newblock High-rank subtensors of high-rank tensors, 2022.
\newblock arXiv:2207.08030.
\newblock \href {https://doi.org/10.48550/ARXIV.2207.08030} {\path{doi:10.48550/ARXIV.2207.08030}}.

\bibitem{KazhdanZ:cohom}
David Kazhdan and Tamar Ziegler.
\newblock Approximate cohomology.
\newblock {\em Selecta Math. (N.S.)}, 24(1):499--509, 2018.
\newblock \href {https://doi.org/10.1007/s00029-017-0335-5} {\path{doi:10.1007/s00029-017-0335-5}}.

\bibitem{KazhdanZ:2020}
David Kazhdan and Tamar Ziegler.
\newblock Properties of high rank subvarieties of affine spaces.
\newblock {\em Geom. Funct. Anal.}, 30(4):1063--1096, 2020.
\newblock \href {https://doi.org/10.1007/s00039-020-00542-4} {\path{doi:10.1007/s00039-020-00542-4}}.

\bibitem{KoppartyMZ}
Swastik Kopparty, Guy Moshkovitz, and Jeroen Zuiddam.
\newblock Geometric rank of tensors and subrank of matrix multiplication.
\newblock In {\em 35th {C}omputational {C}omplexity {C}onference}, volume 169 of {\em LIPIcs. Leibniz Int. Proc. Inform.}, pages Art. No. 35, 21. Schloss Dagstuhl. Leibniz-Zent. Inform., Wadern, 2020.
\newblock \href {https://doi.org/10.4230/LIPIcs.CCC.2020.35} {\path{doi:10.4230/LIPIcs.CCC.2020.35}}.

\bibitem{Landsberg08}
J.~M. Landsberg.
\newblock Geometry and the complexity of matrix multiplication.
\newblock {\em Bull. Amer. Math. Soc. (N.S.)}, 45(2):247--284, 2008.
\newblock \href {https://doi.org/10.1090/S0273-0979-08-01176-2} {\path{doi:10.1090/S0273-0979-08-01176-2}}.

\bibitem{Lovett2019}
Shachar Lovett.
\newblock The analytic rank of tensors and its applications.
\newblock {\em Discrete Anal.}, pages Paper No. 7, 10, 2019.
\newblock \href {https://doi.org/10.19086/da} {\path{doi:10.19086/da}}.

\bibitem{NaslundEGZ}
Eric Naslund.
\newblock Exponential bounds for the {Erd\H{o}s}-{G}inzburg-{Z}iv constant.
\newblock {\em J. Combin. Theory Ser. A}, 174:105185, 19, 2020.
\newblock \href {https://doi.org/10.1016/j.jcta.2019.105185} {\path{doi:10.1016/j.jcta.2019.105185}}.

\bibitem{Naslund2020}
Eric Naslund.
\newblock The partition rank of a tensor and {$k$}-right corners in {$\Bbb F_q^n$}.
\newblock {\em J. Combin. Theory Ser. A}, 174:105190, 25, 2020.
\newblock \href {https://doi.org/10.1016/j.jcta.2019.105190} {\path{doi:10.1016/j.jcta.2019.105190}}.

\bibitem{NaslundS17}
Eric Naslund and Will Sawin.
\newblock Upper bounds for sunflower-free sets.
\newblock {\em Forum Math. Sigma}, 5:Paper No. e15, 10, 2017.
\newblock \href {https://doi.org/10.1017/fms.2017.12} {\path{doi:10.1017/fms.2017.12}}.

\bibitem{ODonnell:2014}
Ryan O'Donnell.
\newblock {\em Analysis of {B}oolean functions}.
\newblock Cambridge University Press, New York, 2014.
\newblock \href {https://doi.org/10.1017/CBO9781139814782} {\path{doi:10.1017/CBO9781139814782}}.

\bibitem{Raz13}
Ran Raz.
\newblock Tensor-rank and lower bounds for arithmetic formulas.
\newblock {\em J. ACM}, 60(6):Art. 40, 15, 2013.
\newblock \href {https://doi.org/10.1145/2535928} {\path{doi:10.1145/2535928}}.

\bibitem{Schechtman2003}
Gideon Schechtman.
\newblock Concentration results and applications.
\newblock In {\em Handbook of the geometry of {B}anach spaces, {V}ol. 2}, pages 1603--1634. North-Holland, Amsterdam, 2003.
\newblock \href {https://doi.org/10.1016/S1874-5849(03)80044-X} {\path{doi:10.1016/S1874-5849(03)80044-X}}.

\bibitem{Schmidt:1985}
Wolfgang~M. Schmidt.
\newblock The density of integer points on homogeneous varieties.
\newblock {\em Acta Math.}, 154(3-4):243--296, 1985.
\newblock \href {https://doi.org/10.1007/BF02392473} {\path{doi:10.1007/BF02392473}}.

\bibitem{Talagrand:1995}
Michel Talagrand.
\newblock Concentration of measure and isoperimetric inequalities in product spaces.
\newblock {\em Inst. Hautes \'{E}tudes Sci. Publ. Math.}, (81):73--205, 1995.

\bibitem{Tao:srank}
Terence Tao.
\newblock \url{https://terrytao.wordpress.com/2016/05/18/a-symmetric-formulation-of-the-croot-lev-pach-ellenberg-gijswijt-capset-bound/}, 2016.

\bibitem{TaoSawin}
Terence Tao and Will Sawin.
\newblock \url{https://terrytao.wordpress.com/2016/08/24/notes-on-the-slice-rank-of-tensors/}, 2016.

\bibitem{TaoZ12}
Terence Tao and Tamar Ziegler.
\newblock The inverse conjecture for the {G}owers norm over finite fields in low characteristic.
\newblock {\em Ann. Comb.}, 16(1):121--188, 2012.
\newblock \href {https://doi.org/10.1007/s00026-011-0124-3} {\path{doi:10.1007/s00026-011-0124-3}}.

\end{thebibliography}
\end{document}
